\chapter*{Preface}
\addcontentsline{toc}{chapter}{Preface}
\markboth{Preface}{Preface}

In August 1900 David Hilbert stood before the Second International Congress of Mathematicians in Paris and set out the problems that he believed would occupy the coming century. The sixth asked for something that sounds almost administrative: a mathematical treatment of the axioms of physics, beginning with probability and mechanics. Behind the modest phrasing sat the deepest unresolved question of nineteenth century science. Boltzmann had shown how the irreversible laws of gases might emerge from reversible collisions of atoms, and his critics had shown, with equal force, that the argument as it stood could not be right. Hilbert asked for the limiting processes, carried out with full rigor, that lead from the atomistic picture to the equations of continua.

A century and a quarter later the problem looks different in three ways. Hilbert's request for axioms of probability was answered by Kolmogorov in 1933~\cite{kolmogorov1933}, and most working scientists no longer remember that it was ever a problem. The mechanics half, after Lanford's theorem of 1975~\cite{lanford1975} and a long program of French and American analysis, has reached the point where Deng, Hani and Ma~\cite{denghanima2024,denghanima2025} have derived the Boltzmann equation from hard spheres for as long as its solution stays regular, and from it the equations of fluids on the torus in dimensions two and three; Chapters~\ref{ch:kinetic} and~\ref{ch:hydro} give the precise statements. And a third reading of the problem has appeared that Hilbert did not pose: the inverse problem. Physics derives continua from particles. Data science must go the other way, inferring initial conditions, forcing and microscopic structure from noisy macroscopic observations, while respecting the physics. The natural language for that inverse direction is Bayesian, and the natural machinery is simulation.

This book is about all three readings. It is written for mathematicians, statisticians, physicists and machine learning researchers who want a single account of how probability, kinetic theory, fluid dynamics and modern Bayesian computation fit together. It is also a companion to \emph{On Hilbert's 8th Problem}~\cite{polsonh8}, a separate volume by one of us on the Riemann hypothesis through Thorin measures and Brownian motion. The two problems are closer than their numbers suggest. Both are, at bottom, problems about the heat equation: the de Bruijn Newman constant is defined by deforming the Riemann $\Xi$ function with the backward heat flow, which multiplies its P\'olya kernel by $e^{tu^2}$, and the Navier Stokes equations are, among other things, a nonlinear heat equation whose backward solution is the central difficulty of inference.

We have given a large place to Henri Poincar\'e. Poincar\'e presided over the same Paris congress at which Hilbert spoke. In the same year he examined Louis Bachelier's thesis on the theory of speculation~\cite{bachelier1900}, an early mathematical treatment of Brownian motion. His recurrence theorem of 1890 was the sharpest objection to Boltzmann's program; his method of arbitrary functions was the first account of how chance emerges from deterministic mechanics; his report in the Dreyfus affair was a public lesson in inverse probability; one of his last papers argued that Planck's quantum hypothesis could not be avoided; and his conjecture of 1904 was solved by Perelman with an entropy functional as the key monotone quantity. A book on the sixth problem that left him out would be missing half its history.

The structure is as follows. Part~\ref{part:problem} states the problem, tells the story of Poincar\'e and Hilbert, sets out the axioms of Kolmogorov and de Finetti and Sanov's large deviation theorem, develops ergodic theory from Boltzmann's hypothesis to Kolmogorov Sinai entropy, and builds the method of types and the theory of moment conditioning, which explain why conditioning a large system on its conserved quantities produces the Maxwellian. Part~\ref{part:brownian} develops the probabilistic toolkit: Brownian motion and Yor's identities read as algorithms, scale mixtures and Thorin measures, entropy and Perelman's functional, and the quantum axioms of von Neumann together with Nelson's stochastic mechanics. Part~\ref{part:fluids} follows the passage from particles to fluids: kinetic theory and propagation of chaos, the Hilbert expansion to Navier Stokes, the Bayesian treatment of Navier Stokes on its solution manifold, which Chapter~\ref{ch:nsinference} gives in summary, and the boundary of that manifold where blowup and turbulence live. Part~\ref{part:manifold} develops that treatment at length, from our paper~\cite{polsonsokolov2026manifold}. It treats the inverse problem of a partial differential equation as Bayesian estimation on the solution manifold. The adjoint evolution of the data is a sufficient statistic for linear equations and, under a uniform Bernstein von Mises statement that we leave open, an asymptotically sufficient one for nonlinear equations. The cost of generative computation depends on the dimension that the data can resolve, and the part works out examples and experiments for the heat, Navier Stokes and Schr\"odinger equations. Part~\ref{part:statmech} turns the picture around and treats inference itself as statistical mechanics: spin systems and phase transitions, generative Bayesian computation as transport, filters and local learning rules, and the thermodynamics of information. A closing chapter lists open problems and builds the bridge to Hilbert's eighth problem.

The book assumes measure theoretic probability and a first course in partial differential equations. The book proves neither the Riemann hypothesis nor the global regularity of the three dimensional Navier Stokes equations, both of which are open, and it says what is known about each (Chapters~\ref{ch:boundary} and~\ref{ch:open}). Of the two theorems that carry the derivational reading, Lanford's is proved in outline and the long time theorem of Deng, Hani and Ma is quoted without proof, with an account of the ideas of its proof (Chapter~\ref{ch:kinetic}).

The chapters depend on one another as follows, apart from passing references. Chapter~\ref{ch:hilbert} also contains a reading guide (Section~\ref{sec:howtoread}).
\begin{itemize}
\item Part~\ref{part:problem}. Chapters~\ref{ch:hilbert} and~\ref{ch:poincare} use no earlier chapter. Chapter~\ref{ch:axioms} uses Chapter~\ref{ch:poincare}, Chapter~\ref{ch:ergodic} uses Chapters~\ref{ch:poincare} and~\ref{ch:axioms}, and Chapter~\ref{ch:types} uses Chapters~\ref{ch:poincare} to~\ref{ch:ergodic}.
\item Part~\ref{part:brownian}. Chapter~\ref{ch:brownian} uses Chapter~\ref{ch:poincare}, Chapter~\ref{ch:mixtures} uses Chapters~\ref{ch:axioms} and~\ref{ch:brownian}, Chapter~\ref{ch:entropy} uses Chapters~\ref{ch:poincare}, \ref{ch:axioms} and~\ref{ch:brownian}, and Chapter~\ref{ch:quantum} uses Chapters~\ref{ch:hilbert} to~\ref{ch:axioms}, \ref{ch:brownian} and~\ref{ch:entropy}.
\item Part~\ref{part:fluids}. Chapter~\ref{ch:kinetic} uses Chapters~\ref{ch:hilbert} to~\ref{ch:brownian}, Chapter~\ref{ch:hydro} uses Chapters~\ref{ch:axioms}, \ref{ch:quantum} and~\ref{ch:kinetic}, Chapter~\ref{ch:nsinference} uses Chapters~\ref{ch:hilbert}, \ref{ch:entropy} and~\ref{ch:hydro}, and Chapter~\ref{ch:boundary} uses Chapters~\ref{ch:mixtures} and~\ref{ch:quantum} to~\ref{ch:nsinference}.
\item Part~\ref{part:manifold}. Chapters~\ref{ch:manifold} to~\ref{ch:examples} form a unit, to be read in order, and are best read after Chapter~\ref{ch:nsinference}, which gives the setting in summary.
\item Part~\ref{part:statmech}. Chapter~\ref{ch:spins} uses Chapters~\ref{ch:poincare}, \ref{ch:mixtures} and~\ref{ch:entropy}, Chapter~\ref{ch:gbc} uses Chapters~\ref{ch:axioms}, \ref{ch:types}, \ref{ch:entropy}, \ref{ch:quantum}, \ref{ch:nsinference} and~\ref{ch:manifold}, Chapter~\ref{ch:filters} uses Chapters~\ref{ch:axioms} to~\ref{ch:types}, \ref{ch:mixtures}, \ref{ch:kinetic} and~\ref{ch:gbc}, Chapter~\ref{ch:kolmogorov} uses Chapters~\ref{ch:axioms}, \ref{ch:gbc} and~\ref{ch:filters}, and Chapter~\ref{ch:open} uses Chapters~\ref{ch:poincare}, \ref{ch:types} to~\ref{ch:entropy}, \ref{ch:nsinference}, \ref{ch:spins} and~\ref{ch:kolmogorov}.
\end{itemize}

The appendices support the text: Appendix~\ref{app:levy} covers L\'evy processes, Thorin measures and related classes, Appendix~\ref{app:code} the code and its reproducibility, Appendix~\ref{app:toolkit} a probability toolkit, Appendix~\ref{app:chronology} a chronology and Appendix~\ref{app:notation} the notation.

Every theorem in the book is proved in the text, proved in an exercise, or quoted from the literature with a reference, and a theorem quoted without its proof says so. The figures were computed or illustrated by code written for this book and for our paper on which Part~\ref{part:manifold} is based.

Much of the material grew out of joint papers written over the last several years: on Bayesian methods for the Navier Stokes equations~\cite{polsonsokolov2026ns}, on estimation on the solution manifold of a partial differential equation~\cite{polsonsokolov2026manifold}, on Yor's identities as algorithms~\cite{polsonsokolovyor}, on Perelman's entropy, dual densities and the Heisenberg uncertainty principle, with Daniel Zantedeschi~\cite{polsonsokolovperelman}, on shot noise processes~\cite{polsonsokolov2026gbcshot}, on generative Bayesian computation~\cite{polsonsokolov2025}, and on gates as approximate nonlinear filters~\cite{polsonsokolov2026gates}. We thank the many colleagues whose comments improved those papers, and we take responsibility for every error that remains.

\medskip
\noindent\textit{Nick Polson, Chicago}\\
\textit{Vadim Sokolov, Fairfax}\\
\textit{October 2026}
